% SHARD-ID: AQM-24-ZARISKI-OPENS-OBSERVABLE-NETS
% SHARD-TITLE: Zariski Opens and Observable Nets
% SHARD-SUMMARY: Defines open-local and closed-complement observable algebras for finite affine spectra.
% SHARD-SUMMARY: Proves that open-local Weyl algebras form a covariant net while closed-complement algebras are contravariant.
% SHARD-SUMMARY: Computes the edge cases \(\operatorname{Spec}\mathbb F_p\) and \(\operatorname{Spec}(\mathbb Z/6\mathbb Z)\).
% SHARD-KEYWORDS: Zariski topology, open set, closed set, observable algebra, Weyl-Heisenberg, Spec, Z/6Z

\section{Zariski Opens and Observable Nets}
\label{sec:zariski-opens-observable-nets}

\subsection{Status and Source Anchors}

This shard stays at the concrete finite-spectrum level.  The Zariski topology
definitions are sourced from the Stacks Project snapshot
\path{references/algebraic_geometry/stacks_project_algebra.tex}: the spectrum,
\(V(T)\), and \(D(f)\) are defined in lines 2972--2983; the Zariski topology
and standard opens are defined in lines 3104--3117.  The finite Weyl local
systems are those of AQM-22 and AQM-23.

The main conclusion is a convention choice.  If the objective is a net whose
inclusions follow inclusions of open sets, then an open \(U\) should carry the
algebra acting nontrivially on \(U\) itself.  The assignment
\(U\mapsto\mathcal A(X\setminus U)\) is also meaningful, but it is a
closed-support or vanishing-locus assignment and its inclusions run in the
opposite direction.

\subsection{Lamport Dependency Skeleton}

\[
\begin{array}{rcl}
D1&:&X=\operatorname{Spec}R,\quad
      V(T)=\{\mathfrak p:T\subset\mathfrak p\},\quad
      D(f)=\{\mathfrak p:f\notin\mathfrak p\}.\\
D2&:&\text{For each }x\in X\text{ choose }
      \mathcal H_x=\ell^2(\kappa(x)),\
      \mathcal A_x=\operatorname{End}(\mathcal H_x).\\
D3&:&E(S)=\bigoplus_{x\in S}\kappa(x)^2,\quad
      \mathcal H(S)=\bigotimes_{x\in S}\mathcal H_x,\quad
      \mathcal A(S)=\bigotimes_{x\in S}\mathcal A_x.\\
L1&:&S\subset T\Rightarrow
      \iota_{S,T}:\mathcal A(S)\hookrightarrow\mathcal A(T)
      \text{ is a functorial unital }*\text{-inclusion}.\\
L2&:&\iota_{S,T}\text{ is the inclusion of Weyl algebras induced by }
      E(S)\hookrightarrow E(T).\\
T1&:&U\mapsto\mathcal A(U)\text{ is the covariant open-local net.}\\
T2&:&U\mapsto\mathcal A(X\setminus U)\text{ is contravariant in opens.}\\
E1&:&X=\operatorname{Spec}\mathbb F_p
      \text{ has only }\emptyset\text{ and }X.\\
E2&:&X=\operatorname{Spec}(\mathbb Z/6\mathbb Z)
      \text{ has two clopen points and algebras }M_2,M_3,M_6.
\end{array}
\]

\subsection{Definitions D1--D3}

Let \(R\) be a commutative ring and let
\[
  X=\operatorname{Spec}R
\]
be the set of prime ideals of \(R\).  For a subset \(T\subset R\), define
\[
  V(T)=\{\mathfrak p\in X:T\subset\mathfrak p\}.
\]
The Zariski closed subsets are the subsets \(V(T)\).  The open subsets are
their complements.  For \(f\in R\), define the standard open
\[
  D(f)=\{\mathfrak p\in X:f\notin\mathfrak p\}=X\setminus V(\{f\}).
\]

In this shard \(X\) is finite.  Fix an ordering of \(X\) only to make tensor
products strict.  For \(x=\mathfrak p\in X\), set
\[
  \mathcal H_x=\ell^2(\kappa(x)),\qquad
  \mathcal A_x=\operatorname{End}_{\mathbb C}(\mathcal H_x),
\]
where \(\kappa(x)\) is the residue field at \(x\).  The local phase label
group is \(\kappa(x)^2\) with the standard alternating form
\[
  \omega_x((q,p),(q',p'))=pq'-qp'.
\]
For any subset \(S\subset X\), define
\[
  E(S)=\bigoplus_{x\in S}\kappa(x)^2,\qquad
  \mathcal H(S)=\bigotimes_{x\in S}\mathcal H_x,\qquad
  \mathcal A(S)=\bigotimes_{x\in S}\mathcal A_x.
\]
The empty direct sum is \(0\), and the empty tensor product is \(\mathbb C\).

\subsection{Functorial Inclusions L1--L2}

\paragraph{Lemma \(L1\).}
If \(S\subset T\subset X\), then
\[
  \iota_{S,T}:\mathcal A(S)\longrightarrow\mathcal A(T),\qquad
  a\longmapsto a\otimes\mathbf 1_{\mathcal H(T\setminus S)}
\]
is an injective unital \(*\)-homomorphism.  If \(S\subset T\subset U\), then
\[
  \iota_{T,U}\circ\iota_{S,T}=\iota_{S,U}.
\]

\emph{Proof.}
The map preserves multiplication, adjoints, and the unit because tensoring
with an identity operator preserves these operations.  It is injective:
if \(a\otimes\mathbf 1=0\), then applying the operator to a simple tensor
\(v\otimes w\) with \(w\ne0\) gives \(av\otimes w=0\) for every \(v\), hence
\(av=0\) for every \(v\), and \(a=0\).  The composition identity is
\[
  (a\otimes\mathbf 1_{T\setminus S})\otimes\mathbf 1_{U\setminus T}
  =
  a\otimes\mathbf 1_{U\setminus S}.
\]

\paragraph{Lemma \(L2\).}
Let \(W_S(e)\) denote the tensor product of local Weyl operators for
\(e\in E(S)\).  Under the extension-by-zero inclusion
\(E(S)\hookrightarrow E(T)\),
\[
  \iota_{S,T}(W_S(e))=W_T(e).
\]
Consequently the finite Weyl-Heisenberg observable algebra generated by
\(E(S)\) is included in the one generated by \(E(T)\).

\emph{Proof.}
For \(x\in S\), the \(x\)-factor of \(W_T(e)\) is the same local Weyl operator
as in \(W_S(e)\).  For \(x\in T\setminus S\), extension by zero gives the
local label \(0\), and \(W_x(0)=\mathbf 1_{\mathcal H_x}\).  Thus
\[
  W_T(e)=W_S(e)\otimes\mathbf 1_{\mathcal H(T\setminus S)}
         =\iota_{S,T}(W_S(e)).
\]
Taking complex spans gives the stated Weyl algebra inclusion.

\subsection{Open-Local Versus Closed-Complement T1--T2}

\paragraph{Theorem \(T1\).}
The assignment
\[
  U\longmapsto\mathcal A_{\rm open}(U):=\mathcal A(U)
\]
from Zariski open subsets \(U\subset X\) to finite-dimensional
\(*\)-algebras is covariant under inclusion: if \(U\subset V\), then
\[
  \mathcal A_{\rm open}(U)\hookrightarrow\mathcal A_{\rm open}(V)
\]
by \(\iota_{U,V}\).  This is the natural choice when one wants observable
algebras with inclusions following open inclusions.

\emph{Proof.}
An inclusion of opens \(U\subset V\) is an inclusion of subsets of the finite
set \(X\).  Lemma \(L1\) gives the unital \(*\)-inclusion, and Lemma \(L2\)
identifies it with the corresponding inclusion of Weyl-Heisenberg algebras.

\paragraph{Theorem \(T2\).}
The assignment
\[
  U\longmapsto\mathcal A_{\rm cl}(U):=\mathcal A(X\setminus U)
\]
acts nontrivially only on the closed complement of \(U\), but it is
contravariant in opens.  If \(U\subset V\), then
\[
  X\setminus V\subset X\setminus U,
\]
so the natural inclusion is
\[
  \mathcal A_{\rm cl}(V)\hookrightarrow\mathcal A_{\rm cl}(U).
\]

\emph{Proof.}
The complement of an open set is closed by definition of open.  Complementing
the inclusion \(U\subset V\) reverses it, giving
\(X\setminus V\subset X\setminus U\).  Lemma \(L1\) then gives the displayed
inclusion.  Thus this construction is valid as a closed-support or
vanishing-locus algebra, but it is not the open-local net of Theorem \(T1\).

\subsection{Edge Case: A Field E1}

Let \(R=\mathbb F_p\).  AQM-22 proves that
\[
  X=\operatorname{Spec}\mathbb F_p=\{\eta\},\qquad \eta=(0).
\]
The only Zariski closed subsets are \(V(0)=X\) and \(V(1)=\emptyset\), hence
the only opens are \(\emptyset\) and \(X\).  The open-local algebras are
\[
  \mathcal A_{\rm open}(\emptyset)=\mathbb C,\qquad
  \mathcal A_{\rm open}(X)=\operatorname{End}(\ell^2(\mathbb F_p))
  \cong M_p(\mathbb C),
\]
with the scalar inclusion \(\lambda\mapsto\lambda\mathbf 1\).  The
closed-complement assignment reverses these:
\[
  \mathcal A_{\rm cl}(\emptyset)=M_p(\mathbb C),\qquad
  \mathcal A_{\rm cl}(X)=\mathbb C.
\]
Thus the one-point field case has no nontrivial intermediate localization;
it only tests the direction of the convention.

\subsection[The Two-Point Case Z/6Z E2]{The Two-Point Case \(\mathbb Z/6\mathbb Z\) E2}

Let \(R=\mathbb Z/6\mathbb Z\).  AQM-23 proves
\[
  X=\operatorname{Spec}R=\{\mathfrak p_2,\mathfrak p_3\},\qquad
  \mathfrak p_2=(2),\quad \mathfrak p_3=(3),
\]
with residue fields \(\kappa(\mathfrak p_2)=\mathbb F_2\) and
\(\kappa(\mathfrak p_3)=\mathbb F_3\).  Since \(2\in\mathfrak p_2\) and
\(2\notin\mathfrak p_3\),
\[
  D(2)=\{\mathfrak p_3\},\qquad V(2)=\{\mathfrak p_2\}.
\]
Since \(3\notin\mathfrak p_2\) and \(3\in\mathfrak p_3\),
\[
  D(3)=\{\mathfrak p_2\},\qquad V(3)=\{\mathfrak p_3\}.
\]
The standard opens therefore contain both singletons, so the topology is
discrete:
\[
  \emptyset,\quad \{\mathfrak p_2\}=D(3),\quad
  \{\mathfrak p_3\}=D(2),\quad X.
\]

With the order \(\mathfrak p_2,\mathfrak p_3\),
\[
  \mathcal H(X)=\mathbb C^2\otimes\mathbb C^3,\qquad
  \mathcal A(X)=M_2(\mathbb C)\otimes M_3(\mathbb C)\cong M_6(\mathbb C).
\]
The open-local and closed-complement algebras are:
\[
\begin{array}{c|c|c}
U & \mathcal A_{\rm open}(U)\subset\mathcal A(X)
  & \mathcal A_{\rm cl}(U)=\mathcal A(X\setminus U)\\
\hline
\emptyset & \mathbb C\,\mathbf 1 & M_2\otimes M_3\\
D(3)=\{\mathfrak p_2\} & M_2\otimes\mathbf 1_3
  & \mathbf 1_2\otimes M_3\\
D(2)=\{\mathfrak p_3\} & \mathbf 1_2\otimes M_3
  & M_2\otimes\mathbf 1_3\\
X & M_2\otimes M_3 & \mathbb C\,\mathbf 1
\end{array}
\]
Thus \(D(3)\), the open where \(3\) does not vanish, carries the qubit algebra
in the open-local convention and the qutrit algebra in the closed-complement
convention.  Likewise \(D(2)\) carries the qutrit algebra open-locally and the
qubit algebra as its closed complement.

Finally,
\[
  [M_2\otimes\mathbf 1_3,\mathbf 1_2\otimes M_3]=0,
\]
and the two singleton open algebras generate
\[
  (M_2\otimes\mathbf 1_3)(\mathbf 1_2\otimes M_3)
  =M_2\otimes M_3.
\]
This is the finite-spectrum analogue of locality: disjoint point supports
commute, and the union of the supports generates the tensor-product algebra.

\subsection{Conclusion}

The best correspondence for the present programme is
\[
  \boxed{U\text{ open}\quad\leadsto\quad
  \mathcal A(U)\text{ acting nontrivially on }U.}
\]
This gives isotony of observable algebras:
\[
  U\subset V\quad\Rightarrow\quad\mathcal A(U)\subset\mathcal A(V).
\]
The alternative
\[
  U\leadsto\mathcal A(X\setminus U)
\]
should be named explicitly as a closed-complement or vanishing-locus
assignment.  It is useful when one wants observables supported where a
function vanishes, but its functorial direction is opposite to open inclusion.
